\paragraph{The cycle and its determinant.}\label{kone:cycle}
Write $A_q=\mathbb Z/q$, $C=\mathbb Z/q$ as an additive group,
$M=2k+3$, $\kappa=1+p^c$ and $a=\kappa^e$, so $e=k+1$.
In the bar complex of $C$ with coefficients in $A_q$ use
\[
 z=\sum_{j_1,\ldots,j_e\in C}[j_1|1|\cdots|j_e|1].
\]
This has degree $2e$. In its boundary the first face sums to zero
because its multiplicity is $q$. The remaining faces cancel in pairs
after shifting a summation variable by one. Put $d_t(x)=(x,tx)$,
$0\le t\le M$. The product-group homology in degree $2e$ with these
coefficients has a generating set of size at most $2e+1=M$ (use the
first projection onto $C$). More explicitly, constants, cyclic
shift minus identity and summation give an exact sequence
from the trivial module through two copies of the cyclic permutation
module to the trivial module. The two connecting sequences and
induction from the kernel bound the generator number in degree $d+2$
by that in degree $d$ plus the kernel-group bounds in degrees $d+1,d+2$.
For a cyclic kernel both bounds are one; degree zero needs one
generator. Hence degree $2e$ needs at most $2e+1$ generators.
The $M+1$ classes $(d_t)_*z$
therefore satisfy a relation $\sum_t n_t(d_t)_*z=\partial B$ with
some $n_{t_0}$ a unit. Indeed their subgroup is a finite abelian
$p$-group with at most $M$ cyclic factors. A redundant generator
modulo $p$ is redundant in the subgroup, by successive lifting
through the powers of $p$. This gives the unit coefficient.

Index the coordinates of $F^{C^2}$ by $j=(j_1,j_2)$.
On coordinate $j$ the diagonal representation $\rho$ of $C^2$
has character $(x,y)\mapsto\zeta^{j_1x+j_2y}$. For $t$ let
$s_t(j)=j_1+tj_2$. The lattice construction uses a finitely supported
integer function $\eta$ on $C^2$. Its sums on the fibers of $s_t$
vanish for $t\ne t_0$. At $t_0$ its pushforward is
\[
 \sum_{i=0}^{M}[X^i]P(X)\,[\kappa^i],\qquad
 P(X)=\prod_{j=0}^{M-1}(X-\kappa^j).
\]
Here $[x]$ denotes a basis vector in $\mathbb Z[C]$. To see the
interpolation, use the product of $[(0,1)]-[(t,0)]$ over $t\ne t_0$
in $\mathbb Z[C^2]$. Each unwanted pushforward kills a factor. At
$t_0$ the product equals a translation times
$\prod_{t\ne t_0}([t_0-t]-1)$. Every $t_0-t$ generates a subgroup
containing $p^c$ in $C$, since $0<|t_0-t|\le M<p^c$. Thus each factor
divides $[p^c]-1$. The desired pushforward is divisible by
$([p^c]-1)^M$: in polynomial coordinates apply
$\prod_{j<M}(U-\kappa^j)$ to $X$, where $Uf(X)=f(X^\kappa)$.
The operator $U-\kappa^j$ carries $(X^{p^c}-1)^j$ times any polynomial
into $(X^{p^c}-1)^{j+1}$ times a polynomial, by factoring
$X^{\kappa p^c}-1$ and reducing the quotient modulo $X^{p^c}-1$.
Lift the quotient in the group ring by $x\mapsto(x,0)$ and undo the
translation. This constructs $\eta$ with both specified properties.

Modify the standard lattice coordinate $j$ by the fractional ideal
with divisor $\eta(j)\mathfrak d$. For each $t$ construct an
$L'$-linear $g_t$ taking the standard lattice to this modified lattice
and commuting with $\rho d_t$. The determinant-one identifications on
blocks use the Steinitz relation for a direct sum of fractional ideals:
one can move the sum of two divisors to one coordinate, and iterate.
Scaling one distinguished coordinate realizes a prescribed determinant
when its divisor is the difference of the two total divisors.
For $t\ne t_0$ all block determinants can be taken as one. At $t_0$
the construction groups coordinates in Galois orbits and leaves the
determinant
\[
 \beta=\prod_i\gamma^{-i}(\alpha)^{[X^i]P_e(X)},\qquad
 P_e(X)=\prod_{0\le j<M,\ j\ne e}(X-\kappa^j)
\]
at character $1$. Indeed the divisor of $\gamma^{-i}\alpha$ at $P'$
is $\mathfrak d(\gamma^{i+1}P')-a\mathfrak d(\gamma^iP')$.
Multiplying this difference by the coefficients of $P_e$ gives
precisely the divisor prescribed by $P=(X-a)P_e$. This explains
both the inverse action in $\beta$ and the missing factor in $P_e$.

Let $H$ be the stabilizer of the standard lattice inside the localized
group containing all $g_t$ and $\rho(C^2)$. One base-field denominator
excludes the support of the divisors and the bad primes of the splitting
map. Set $\rho'=g_{t_0}^{-1}\rho g_{t_0}$ and $h_t=g_{t_0}^{-1}g_t$.
Both representations take values in $H$, and $h_t\in H$ since both
maps $g_t,g_{t_0}$ take the standard lattice to the same lattice.
For the bar prism $P_h$ one has
$\partial P_h+P_h\partial=(\mathrm{conj}_{h^{-1}})_*-1$.
Consequently the chain
\[
 Y=\rho'_*B-\rho_*B+\sum_t n_t P_{h_t}((\rho'd_t)_*z)
\]
is closed in $H$. In fact $h_t^{-1}\rho'd_th_t=\rho d_t$, so the
boundary of the sum is $\rho_*\partial B-\rho'_*\partial B$.
In the localized group $Y$ is homologous to
$\sum_t n_t P_{g_t}((\rho d_t)_*z)$. This follows by applying the
prism identity to $\rho_*B$ and the prism product identity to
$g_t=g_{t_0}h_t$ on each cycle. Multiply $Y$ by $n_{t_0}^{-1}$.
The standard lattice is finite projective over $\mathcal O_{L'}$.
Adding a projective complement with trivial action realizes this cycle
in some $\mathrm{GL}_N(\mathcal O_{L'})$.

At an allowed residue field $k_0$, restriction of scalars splits into
the conjugate blocks indexed by $\sigma\in\mathrm{Gal}(F/L')$.
For a character $v:k_0^\times\to A_q$, the finite-field functional
on a scalar-character block of weight $x$ evaluates a prism for $g$
as $x^e v(\det g)$. On the scalar line this is the product of $e$
cyclic carry cocycles with $v$: each consecutive pair in $z$ sums to
$x$, and insertion in the prism leaves $v(g)$. Transfer over projective
lines gives the additive functional on vector spaces. The line count
for direct sums uses $|k_0|=1$ in $A_q$, since the reduced $\zeta$ has
order $q$. Additivity kills the trivial complement. Blocks of determinant
one give zero. In the surviving block let $\sigma\zeta=\zeta^{\chi(\sigma)}$.
Summing the conjugate blocks gives
\[
 \sum_\sigma \chi(\sigma)^e v(\overline{\sigma\beta})
   =P_e(a)\,v\big(\overline{\operatorname{nrm}(\alpha)}\big).
\]
Indeed replacing $\alpha$ by $\gamma^{-i}\alpha$ multiplies the
weighted sum by $a^i$ after reindexing, and the sum for $\alpha$ is
the evaluation of its twisted norm. The unit coefficient cancels the
normalization of $Y$. For a cyclic residue unit group take $v$ with
kernel the $q$th powers. Each factor $\kappa^e-\kappa^j$ of $P_e(a)$
has valuation $c+v_p(|e-j|)\le c+M$, by lifting the exponent
(with $c\ge2$ at $p=2$). In $A_q$, $P_e(a)$ therefore divides
$p^{M(c+M)}$. Equality of two evaluations scaled by $p^{2e}$
implies that the ratio of their twisted norms raised to
$p^{2e+M(c+M)}=p^b$ is a $q$th power at this prime.

\paragraph{Detecting powers and bounding the image.}\label{kone:rank}
If $\xi\in F^\times$ passes all these residue tests, choose
$x$ with $x^q=\xi$ in an algebraic closure. For each automorphism
$g$ over $F$, the condition $\sigma x=gx$ describes an open coset.
Degree-one Chebotarev gives a Frobenius in this coset outside the finite
exclusion set and the primes dividing $q$. Its residue field has prime
order, so reduction from $\mathcal O_{L'}$ is surjective. In the residue
field $\xi$ has a $q$th root and all the $q$th roots of unity are
present. Hence the residue of $x$ is fixed by Frobenius. The quotient
$\sigma x/x$ is a $q$th root of unity congruent to one. If it were
different from one, its geometric sum of length $q$ would be zero,
giving $q=0$ in the residue field. Thus $gx=x$ for every $g$ and
$x\in F$. Applying this argument to the ratio proves injectivity
for the simultaneous evaluations of the finite set.

For the image bound put $m_1=2e$, and stabilize the finitely many
cycles to rank $n\ge2(m_1+1)+4$. An integral lift $Z$ of a
mod-$q$ cycle has $\partial Z=qV$. The homology class of $V$ is
$q$-torsion in degree $m_1$, independent of the lift. When this class
vanishes one can subtract $q$ times an integral chain from $Z$ to
make it a cycle. There are at most $T_{m_1}$ fibers of this obstruction
by homological stability. Differences in a fiber therefore reduce the
count to evaluations on integral cycles.

Write $\psi_p$ for the $p$-fold block sum on integral homology and
$i$ for stabilization to the same rank. Starting with $z_0=x$ put
$z_{t+1}=(\psi_p)_*z_t-p^{t+2}i_*z_t$. A primitive evaluation
$\phi$ is stable and additive on block sums. Hence
$\phi(z_{t+1})=p(1-p^{t+1})\phi(z_t)$. The factors $1-p^{t+1}$
are units modulo $q$. The scaled evaluations $p^{m_1}\phi(x)$ thus
factor through the integral span of the $z_{m_1}$. Its rank is at
most $r_2(L')$. To obtain this last inequality the source passes to
real cycles and pairs with one odd-degree functional at each infinite
place. A class in the common kernel, if nonzero, survives stabilization
and is separated by a closed cochain. The surjectivity theorem represents
that cochain on cycles by an invariant form. Its weight decomposition
consists of the place functionals and terms of weights $j=2,\ldots,m_1+1$.
The latter vanish on the displayed span: the factor at $t=j-2$ is
$p^j-p^j=0$. The former vanish by the kernel condition. This contradicts
separation. The real image thus has dimension at most the number of
complex places, giving the integral rank bound.

For any finite set of evaluations use the finitely generated subgroup
spanned by lifts of its elements. Its free part contributes at most
$q^{r_2(L')}$ values. Its torsion contributes at most $T_{m_1+1}$
by stability, also at the enlarged rank for $z_{m_1}$. Multiplication
by the number of obstruction fibers gives the claimed bound.

\paragraph{Transfer over residue-field lines.}\label{kone:lines}
Here is the transfer used in the cycle calculation of
Paragraph~\ref{kone:cycle}. Let $V$ be a vector space over the finite
residue field $k_0$, let $\rho:G\to\mathrm{GL}(V)$, and let $t$ be the
functional on bars of $k_0^\times$ which vanishes on boundaries.
Choose a representative on each line. This gives a function
$\lambda:V\setminus\{0\}\to k_0^\times$ with
$\lambda(sv)=s\lambda(v)$, by writing $v$ as a multiple of its
representative. Give the lines the right action
$L\cdot g=\rho(g)^{-1}L$. For $0\ne v\in L$, put
\[
 c(g,L)=\lambda(v)/\lambda(\rho(g)^{-1}v).
\]
The ratio is independent of $v$. It satisfies
$c(gh,L)=c(g,L)c(h,L\cdot g)$ by cancellation at the intermediate
line. For a bar $[g_1|\cdots|g_n]$ sum, over initial lines $L_0$,
the bars $[c(g_1,L_0)|\cdots|c(g_n,L_{n-1})]$, where
$L_i=L_{i-1}\cdot g_i$. This sum commutes with the bar boundary.
For the first face reindex by $L_1$; for an interior face use the
cocycle identity; the last face just forgets the last arrow.
Composition with $t$ therefore kills boundaries.

If $\lambda$ is multiplied by a scalar $b(L)$ on line $L$, each
edge label is multiplied by $b(L_{i-1})/b(L_i)$. The two labelled
paths are related by the vertical edges $b(L_i)$: the products along
the two sides of each square agree. Triangulate the prism on each
path by inserting the vertical edge successively at the $n+1$
vertices, with alternating signs. Its interior faces cancel in pairs;
the two ends give the difference of the old and new bars, and the
side faces are the corresponding prisms on the boundary. Thus the
two sums agree under $t$ on cycles. In particular a linear change
of coordinates, which transports both lines and representatives,
preserves the resulting functional.

On $V\oplus W$, send a line to its projection in $V$ when that
projection is nonzero, and otherwise to its line in $W$. Use the
normalizing function of the indicated projection. Each line above
a fixed line in $V$ is the graph of a linear map from that line to
$W$, so its fiber has $|k_0|^{\dim W}$ elements. A line in $W$ has
one preimage. For block diagonal actions the edge labels in a fiber
are exactly those of its projected path. Both fiber sizes equal one
as coefficients in $\mathbb Z/q$, since $q\mid |k_0|-1$. The
functional on a direct sum is consequently the sum of the two
functionals on cycles. On the one-dimensional scalar representation
there is a single line and $c(g,L)=g$, so the construction recovers
$t$. These two facts justify passage from the scalar computation to
the sum over the residue-field blocks.

